% !TeX program = pdflatex
% =============================================================================
% Arbeitsblatt 3: Kraftaddition & Kraftzerlegung
% UZH Praktikum I -- Sara Romer Urban, Kantonsschule im Lee, HS2026
% Klasse: 2e (02.09.2026)
%
% Kompakte Fassung nach Ueberarbeitung durch Loris (01.-02.09.2026): keine
% Lernziele-Box, keine Einleitung -- startet direkt mit der vergroesserten
% Feder-Masse-Skizze (samt Koordinatensystem/Kraeftediagramm-Inset) und
% einem kurzen Kraeftegleichgewicht-Satz statt der fruaeheren Zeichen-
% aufgabe. Die Vektor-Theorie nutzt das Seilziehen-Beispiel
% (Bilder/seilziehen.png) als Illustration, dass Betrag \emph{und} Richtung
% zaehlen. Kraefteparallelogramm- und Kraftzerlegung-Theorie verwenden
% fertige Illustrationen aus Bilder/ statt eigener tikz-Diagramme
% (kraefte_parallelogramm.png, lampe_aufgabe.png). Neu: Aufgabe zur
% Addition von fuenf Kraeften (Bilder/5kraefte_addition_aufgabe(.png/
% _loesung.png)). Die fruehere Kernaufgabe (Winkelabhaengigkeit der
% Resultierenden), das Waescheseil-Beispiel und die freie
% Alltagsbeispiel-Aufgabe wurden gestrichen; stattdessen schliesst das
% Blatt mit einer Partnerarbeit, in der die SuS selbst eine Situation
% erfinden und einander zum Loesen weitergeben. Reine Uebungsaufgaben
% (Schiff, Seilbahn, Hundeschlitten, Seilzug mit drei Gewichten) stehen im
% separaten aufgabenblatt2-kraefte-addieren-zerlegen.tex.
%
% Quellen: lernziele.md (dieser Ordner, Planungsstand vor der
% Ueberarbeitung -- weicht in Aufbau/Umfang inzwischen von diesem Blatt ab),
% kraefte-addieren-zerlegen-lektionsplan.tex (dieser Ordner),
% arbeitsblatt2-gewichtskraft-federkraft.tex (Feder-Masse-Skizze,
% Vorlage fuer Box-Definitionen ohne Farbrahmen).
% =============================================================================
\documentclass[addpoints,11pt,a4paper]{exam}

\usepackage[T1]{fontenc}
\usepackage[utf8]{inputenc}
\usepackage[ngerman]{babel}
\usepackage{lmodern}
\usepackage[a4paper,margin=2.2cm,headheight=14pt]{geometry}
\usepackage{amsmath}
\usepackage{siunitx}
% Hausstil: zusammengesetzte Einheiten immer als echter Bruch (\frac), nicht
% als Inline-Schraegstrich oder \tfrac -- gilt global fuer jedes \si/\SI.
\sisetup{per-mode=fraction,fraction-command=\frac}
\usepackage{array,booktabs,tabularx,multirow}
\usepackage{enumitem}
\usepackage{xcolor}
\usepackage{colortbl}
\usepackage{tikz}
\usepackage{physics}
\usepackage[most]{tcolorbox}
\usepackage{lastpage}
\usepackage{graphicx}
\usepackage{hyperref}

\usetikzlibrary{arrows.meta,calc,angles,quotes,decorations.pathmorphing}

\usepackage{setspace}
\usepackage{needspace}
\setlength{\parindent}{0pt}
\setlength{\parskip}{0.6em}
\setstretch{1.08}
\setlist[itemize]{itemsep=0.4em}

% -----------------------------
% Farben (Corporate-Farbschema Physik KFR)
% -----------------------------
\definecolor{titleblue}{HTML}{183A6B}
\definecolor{accentblue}{HTML}{2E6F9E}
\definecolor{lightblue}{HTML}{EAF4FB}
\definecolor{lightgray}{HTML}{F4F6F8}
\definecolor{darkgray}{HTML}{4A4A4A}
\definecolor{accentorange}{HTML}{D9720C}
\definecolor{accentpurple}{HTML}{7B2D8E}
% Hinweisbox: Orange (accentorange), fuer wichtige Klarstellungen.
\definecolor{lightorange}{HTML}{FCEEE0}
% Formelbox: Violett (accentpurple), fuer die zentralen Merksaetze.
\definecolor{lightpurple}{HTML}{F3EAF7}

% Links (hyperref): farblich ans Corporate-Farbschema angeglichen.
\hypersetup{
  colorlinks=true,
  urlcolor=accentblue,
  linkcolor=accentblue,
  pdftitle={Arbeitsblatt 3: Kraftaddition \& Kraftzerlegung},
  pdfauthor={Loris Keller}
}

% -----------------------------
% Spaltentypen
% -----------------------------
\newcolumntype{Y}{>{\centering\arraybackslash}X}
\newcolumntype{P}[1]{>{\raggedright\arraybackslash}p{#1}}

% -----------------------------
% Boxen
% -----------------------------
\tcbset{before skip=10pt, after skip=10pt}
\newtcolorbox{titlebox}{
  enhanced,
  colback=titleblue,
  colframe=titleblue,
  boxrule=0pt,
  arc=3mm,
  left=3mm,right=3mm,top=3mm,bottom=3mm
}

\newlength{\aufgabenindent}
\settowidth{\aufgabenindent}{10.\hskip\labelsep}

% Praktikum I: Lernziele/Einleitung/Theorie/Aufgaben werden NICHT mehr
% farbig gerahmt (nur Formel-, Hinweis- und Antwortbox bleiben Boxen) --
% siehe institutionen/UZH-Praktikum-I/CLAUDE.md, Abschnitt "Boxen-Layout".
\newenvironment{infobox}{%
  \par\addvspace{10pt}\noindent
}{%
  \par\addvspace{10pt}
}

\newenvironment{theoriebox}[1][Theorie]{%
  \par\addvspace{10pt}\noindent\textbf{#1}\par\smallskip\noindent
}{%
  \par\addvspace{10pt}
}

% taskbox/groupbox stehen in questions VOR dem ersten \question (Bilder,
% Fliesstext) -- exam.cls' questions-Liste braucht dafuer eine echte,
% eigenstaendige Box (sonst "Something's wrong--perhaps a missing \item",
% weil z.B. ein \begin{center} vor dem ersten \item der Aussenliste steht).
% Deshalb bleiben es tcolorboxen, nur unsichtbar gemacht (keine Farbe/Rahmen,
% kein Padding) statt sie durch reine Absaetze zu ersetzen.
\newtcolorbox{taskbox}[1][]{
  enhanced, breakable,
  colback=white, colframe=white, boxrule=0pt,
  left=0mm,right=0mm,top=0mm,bottom=0mm,
  before skip=10pt, after skip=10pt,
  before upper={%
    \def\tmpboxtitle{#1}%
    \ifx\tmpboxtitle\empty\else\noindent\textbf{#1}\par\smallskip\fi
  }
}

\newtcolorbox{groupbox}[1][Gruppenarbeit]{
  enhanced, breakable,
  colback=white, colframe=white, boxrule=0pt,
  left=0mm,right=0mm,top=0mm,bottom=0mm,
  before skip=10pt, after skip=10pt,
  before upper={\noindent\textbf{#1}\par\smallskip}
}

% beispielbox: fuer die Einleitung -- wie die anderen Inhaltsboxen ungerahmt.
\newenvironment{beispielbox}[1][Beispiel]{%
  \par\addvspace{10pt}\noindent\textbf{#1}\par\smallskip\noindent
}{%
  \par\addvspace{10pt}
}

% hinweisbox: Orange, fuer wichtige Klarstellungen/Verwechslungswarnungen.
\newtcolorbox{hinweisbox}[1][Hinweis]{
  enhanced,
  breakable,
  colback=lightorange,
  colframe=accentorange,
  title={#1},
  fonttitle=\bfseries,
  boxrule=0.7pt,
  arc=2mm,
  left=5mm,right=5mm,top=2mm,bottom=2mm
}

% formelbox: Violett, um die zentralen Merksaetze dieser Einheit hervorzuheben.
\newtcolorbox{formelbox}[1][Merksatz]{
  enhanced,
  breakable,
  colback=lightpurple,
  colframe=accentpurple,
  title={#1},
  fonttitle=\bfseries,
  boxrule=0.9pt,
  arc=2mm,
  left=5mm,right=5mm,top=2mm,bottom=2mm
}

\newtcolorbox{answerbox}{
  breakable,
  colback=white,
  colframe=gray!55,
  boxrule=0.5pt,
  arc=1.5mm,
  left=2mm,right=2mm,top=2mm,bottom=2mm
}

% Marke: kleines farbiges Inline-Label fuer Teilschritte/Ergebnis-Callouts.
\newtcbox{\Marke}[1][accentpurple]{
  nobeforeafter,
  colback=#1,
  colframe=#1,
  coltext=white,
  boxrule=0pt,
  arc=2pt,
  left=5pt,right=5pt,top=2pt,bottom=2pt,
  fontupper=\bfseries\sffamily\small
}

% --- Feinraster-Helfer fuer alle grafischen Konstruktionsaufgaben ----------
% Feld mit feinem Punkt-/Liniengitter (0.5 Kaestchen Nebenlinien, 1 Kaestchen
% Hauptlinien), sodass bei Massstab 1 Kaestchen = 1 N jede Kraft exakt und
% nachzaehlbar eingezeichnet werden kann. #1#2 = linke untere Ecke, #3#4 =
% rechte obere Ecke.
\newcommand{\Feinraster}[4]{%
  \draw[gray!22,step=0.5] (#1,#2) grid (#3,#4);
  \draw[gray!55,step=1] (#1,#2) grid (#3,#4);
}

% Bild-Helfer: zentriertes Bild mit spuerbarem Abstand zum umgebenden Text
% (statt nur des normalen Absatzabstands). #1 = Breite, #2 = Dateipfad.
\newcommand{\Bild}[2]{%
  \par\vspace{0.6cm}
  \begin{center}
    \includegraphics[width=#1]{#2}
  \end{center}
  \vspace{0.6cm}
}

% --- Loesungen ein-/ausblenden -----------------------------------------------
\noprintanswers
%\printanswers

\pointformat{}
\renewcommand{\solutiontitle}{\noindent\color{titleblue}\textbf{L\"osung:}\enspace}

\pagestyle{headandfoot}
\cfoot{\textcolor{darkgray}{Seite \thepage\ von \pageref{LastPage}}}

\newcommand{\Kopfzeile}[4]{%
  \begin{titlebox}
    {\color{white}\sffamily\bfseries\Large #4}
  \end{titlebox}
  \vspace{0.5em}
  \noindent\textbf{Klasse:}~#1 \hfill \textbf{Datum:}~#2 \hfill \textbf{Thema:}~#3
    \hfill \textbf{Lehrperson:}~Loris Keller
  \vspace{0.8em}
  \lhead{\textcolor{darkgray}{#3}}
  \rhead{\textcolor{darkgray}{#4}}
}

\newlength{\antwortraumlen}
\newcommand{\Antwortraum}[1]{%
  \pgfmathsetlength{\antwortraumlen}{#1*2.5cm}%
  \begin{answerbox}
    \rule{0pt}{\antwortraumlen}%
  \end{answerbox}%
}

% Werte fuer Kopfzeile zentral setzen
\newcommand{\Klasse}{2e}
\newcommand{\Datum}{02.09.2026}
\newcommand{\Thema}{Kraftaddition \& Kraftzerlegung}
\newcommand{\ArbeitsblattTitel}{Arbeitsblatt 3: Kraftaddition \& Kraftzerlegung}

\begin{document}

\Kopfzeile{\Klasse}{\Datum}{\Thema}{\ArbeitsblattTitel}

% =============================================================================
% Repetition: Feder-Masse-Skizze mit Koordinatensystem (LZ1)
% =============================================================================
\textbf{Repetition}

\begin{center}
\begin{tikzpicture}[>=Latex,scale=1.8]
  \draw[very thick] (-1,3) -- (2,3);
  \foreach \x in {-0.8,-0.4,...,1.8} { \draw (\x,3) -- (\x-0.25,3.3); }
  \draw[decorate,decoration={coil,aspect=0.3,segment length=3.2pt,amplitude=3pt}] (0.5,3) -- (0.5,1.6);
  \draw[fill=lightgray!60] (0.1,1.1) rectangle (0.9,1.6);
  \node at (0.5,1.35) {$m$};
  \fill (0.5,1.35) circle (0.5pt);

  \coordinate (O) at (3.8,1.75);
  \draw[gray,dashed] (O) -- ++(0,-1.15);
  \draw[->,gray] (O) -- ++(1.15,0) node[right]{\scriptsize $x$};
  \draw[->,gray] (O) -- ++(0,1.15) node[above]{\scriptsize $y$};
  \draw[->,ultra thick,titleblue] (O) -- ++(0,1) node[midway,right]{$\vec F$};
  \draw[->,ultra thick,accentorange] (O) -- ++(0,-1) node[midway,right]{$\vec F_G$};
  \fill (O) circle (1pt);
\end{tikzpicture}
\end{center}
{\scriptsize Links: Feder-Masse-System. Rechts: dieselben Kräfte im
Koordinatensystem.}

Im Gleichgewicht gilt für die Vektorsumme aller Kräfte:
$\vec F_{\text{res}} = \vec F + \vec F_G = \vec 0$.

% =============================================================================
% Theorie: Kraft als vektorielle Groesse
% =============================================================================
\begin{theoriebox}[Kraft als vektorielle Grösse]
Eine Kraft ist eine \textbf{\emph{vektorielle Grösse}}: Sie hat einen Betrag
und eine Richtung und wird deshalb immer in einem
\textbf{\emph{Koordinatensystem}} beschrieben. In der Skizze oben zeigt die
Gewichtskraft $\vec F_G$ in die negative $y$-Richtung, die Federkraft
$\vec F$ in die positive $y$-Richtung. Es gibt aber auch Kräfte, die
gleichzeitig eine $x$- und eine $y$-Komponente besitzen, also schräg im
Koordinatensystem liegen.

Ein einfaches Beispiel dafür, dass die Richtung einer Kraft genauso wichtig
ist wie ihr Betrag: Beim Tauziehen ziehen zwei Mannschaften in
entgegengesetzte Richtungen. Welche Mannschaft gewinnt, hängt vom Betrag
\emph{beider} Kräfte ab, nicht nur vom eigenen.

\Bild{15.5cm}{Bilder/seilziehen.png}

\begin{hinweisbox}[Hinweis: Kraftmassstab]
Zeichnet man mehrere Kraftpfeile in dasselbe Diagramm, muss für alle
derselbe Kraftmassstab gelten (z.\,B. $1$ Kästchen $\hat=$ $\SI{1}{\newton}$).
Sonst lässt sich das Ergebnis nachher nicht mehr nachmessen.
\end{hinweisbox}
\end{theoriebox}

% =============================================================================
% Theorie: Kraftaddition zweier Kraefte, Kraefteparallelogramm
% =============================================================================
\begin{theoriebox}[Kraftaddition zweier Kräfte: das Kräfteparallelogramm]
Wirken zwei Kräfte $\vec F_1$ und $\vec F_2$ am \textbf{\emph{selben
Angriffspunkt}} eines Körpers, lässt sich ihre gemeinsame Wirkung durch eine
einzige Ersatzkraft beschreiben, die \textbf{\emph{Resultierende}}
$\vec F_{\text{res}}$. Sie wird mit dem \textbf{\emph{Kräfteparallelogramm}}
konstruiert:

\begin{enumerate}[leftmargin=1.2cm,itemsep=0.2em]
  \item Beide Kräfte massstabsgetreu vom gemeinsamen Angriffspunkt aus
  einzeichnen.
  \item Das Parallelogramm ergänzen: An die Spitze von $\vec F_1$ eine Parallele
  zu $\vec F_2$ zeichnen, und an die Spitze von $\vec F_2$ eine Parallele zu
  $\vec F_1$.
  \item Die Resultierende $\vec F_{\text{res}}$ ist die \textbf{\emph{Diagonale}}
  des Parallelogramms, die durch den gemeinsamen Angriffspunkt verläuft (nicht
  die andere Diagonale!).
\end{enumerate}

\Bild{15.5cm}{Bilder/kraefte_parallelogramm.png}
{\scriptsize Links: die beiden Kräfte $\vec F_1$, $\vec F_2$ vom gemeinsamen
Angriffspunkt aus. Rechts: das ergänzte Parallelogramm mit der Resultierenden
$\vec F_{\text{ges}}$ als Diagonale.}
\end{theoriebox}

% =============================================================================
% Uebung: Fuenf Kraefte addieren
% =============================================================================
\needspace{6cm}
\begin{questions}
\begin{taskbox}[Aufgabe: Fünf Kräfte addieren]
\question[6] Am selben Angriffspunkt greifen fünf Kräfte $\vec F_1$ bis
$\vec F_5$ in unterschiedliche Richtungen an:

\Bild{13cm}{Bilder/5kraefte_addition_aufgabe.png}

Konstruieren Sie mit dem Kopf-an-Schwanz-Verfahren die Resultierende
$\vec F_{\text{Ges}}$ aller fünf Kräfte: Hängen Sie die Pfeile (in
beliebiger Reihenfolge) mit gleichbleibender Länge und Richtung
aneinander, und zeichnen Sie den Pfeil vom ersten Schwanz zur letzten
Spitze ein.

\begin{center}
\begin{tikzpicture}[>=Latex,scale=1]
  \Feinraster{0}{0}{10}{8}
\end{tikzpicture}
\end{center}

\begin{solution}
\Bild{12.5cm}{Bilder/5kraefte_addition_aufgabe_loesung.png}
\end{solution}
\end{taskbox}
\end{questions}

% =============================================================================
% Theorie: Kraftzerlegung mit dem Kraefteparallelogramm
% =============================================================================
\begin{theoriebox}[Kraftzerlegung: die Umkehrung der Kraftaddition]
Manchmal ist der umgekehrte Fall gefragt: Eine Kraft ist gegeben, und man
möchte wissen, welche zwei Teilkräfte in zwei \emph{vorgegebenen} Richtungen
zusammen dieselbe Wirkung hätten. Das ist die \textbf{\emph{Kraftzerlegung}},
dieselbe Konstruktion wie bei der Kraftaddition, nur in umgekehrter
Reihenfolge:

\begin{enumerate}[leftmargin=1.2cm,itemsep=0.2em]
  \item Die gegebene Kraft $\vec F$ sowie die beiden vorgegebenen Richtungen
  vom selben Angriffspunkt aus einzeichnen.
  \item Von der Spitze von $\vec F$ aus Parallelen zu beiden Richtungen
  ziehen, bis sie diese schneiden.
  \item Die beiden Teilkräfte $\vec F_1$ und $\vec F_2$ sind die Strecken vom
  Angriffspunkt bis zu den beiden Schnittpunkten.
\end{enumerate}

\begin{hinweisbox}[Hinweis: Addition oder Zerlegung?]
Bei der \textbf{\emph{Addition}} sind die Beträge und Richtungen \emph{aller}
Kräfte bereits bekannt, gesucht ist die Resultierende. Bei der
\textbf{\emph{Zerlegung}} ist es umgekehrt: Nur \emph{eine} Kraft und die
beiden \emph{Richtungen} der gesuchten Teilkräfte sind bekannt, ihre Beträge
sind gesucht.
\end{hinweisbox}

\textbf{Beispiel:} Eine Lampe hängt an zwei Seilen zwischen zwei Wänden. Die
Gewichtskraft $\vec F_G$ wird in die beiden Seilrichtungen zerlegt.

\Bild{15.5cm}{Bilder/lampe_aufgabe.png}
\end{theoriebox}

% =============================================================================
% Partnerarbeit: eigene Situation erfinden (Addition oder Zerlegung)
% =============================================================================
\begin{questions}
\begin{groupbox}[Partnerarbeit: Erfinden Sie Ihre eigene Situation]
\question Denken Sie sich zu zweit eine eigene Situation aus, in der
mehrere Kräfte auf einen Körper wirken (Addition) oder eine Kraft in zwei
Richtungen zerlegt werden kann (Zerlegung). Seien Sie kreativ: Die Situation
darf auch ungewöhnlich oder erfunden sein, muss aber physikalisch sinnvoll
sein.

\begin{parts}
  \part[3] Skizzieren Sie Ihre Situation unten: Zeichnen Sie den Körper, den
  Angriffspunkt sowie entweder mehrere Kraftpfeile (Addition) oder eine Kraft
  mit zwei vorgegebenen Richtungen (Zerlegung) massstabsgetreu ein. Geben Sie
  den verwendeten Kraftmassstab an, und bei einer Zerlegung zusätzlich die
  beiden Richtungen, falls diese nicht bereits aus der Skizze ersichtlich
  sind.
  \Antwortraum{2.5}

  \part[5] Tauschen Sie die Skizze mit Ihrem Partner/Ihrer Partnerin.
  Bestimmen Sie grafisch (je nachdem, was Ihnen vorgegeben wurde) entweder
  die Resultierende oder die beiden Teilkräfte, und tragen Sie die
  Konstruktion direkt in die Skizze oben ein.
\end{parts}
\end{groupbox}
\end{questions}

\end{document}
